% theory_cy2_cy3_fibration.tex
% Research note: The CY_2 -> CY_3 fibration and the Borcherds lift
% Raeez Lorgat, 2 April 2026

\documentclass[11pt]{article}
\usepackage[T1]{fontenc}
\usepackage[utf8]{inputenc}
\usepackage{lmodern}
\frenchspacing
\usepackage{amsmath,amssymb,amsthm}
\usepackage{mathtools}
\usepackage{enumitem}
\usepackage{geometry}
\geometry{margin=1.2in}
\usepackage[unicode]{hyperref}

% Theorem environments
\theoremstyle{plain}
\newtheorem{theorem}{Theorem}[section]
\newtheorem{proposition}[theorem]{Proposition}
\newtheorem{lemma}[theorem]{Lemma}
\newtheorem{corollary}[theorem]{Corollary}
\newtheorem{conjecture}[theorem]{Conjecture}
\theoremstyle{definition}
\newtheorem{definition}[theorem]{Definition}
\newtheorem{construction}[theorem]{Construction}
\newtheorem{example}[theorem]{Example}
\theoremstyle{remark}
\newtheorem{remark}[theorem]{Remark}
\newtheorem{warning}[theorem]{Warning}

% Macros (matching main monograph)
\newcommand{\CY}{\mathrm{CY}}
\newcommand{\HH}{\mathrm{HH}}
\newcommand{\HC}{\mathrm{HC}}
\newcommand{\Tr}{\mathrm{Tr}}
\newcommand{\Ran}{\mathrm{Ran}}
\newcommand{\Rep}{\mathrm{Rep}}
\newcommand{\Coh}{\mathrm{Coh}}
\newcommand{\Fuk}{\mathrm{Fuk}}
\newcommand{\Mod}{\mathrm{Mod}}
\newcommand{\Sym}{\mathrm{Sym}}
\newcommand{\End}{\mathrm{End}}
\newcommand{\Hom}{\mathrm{Hom}}
\newcommand{\id}{\mathrm{id}}
\newcommand{\op}{\mathrm{op}}
\newcommand{\Eone}{E_1}
\newcommand{\Etwo}{E_2}
\newcommand{\Einf}{E_\infty}
\newcommand{\frakg}{\mathfrak{g}}
\newcommand{\cC}{\mathcal{C}}
\newcommand{\cA}{\mathcal{A}}
\newcommand{\cR}{\mathcal{R}}

\title{The $\CY_2 \to \CY_3$ Fibration and the Borcherds Lift\\[6pt]
\large\itshape From Jacobi Forms to Siegel Modular Forms\\
via Elliptic Fibration of Quantum Vertex Chiral Groups}
\author{Raeez Lorgat}
\date{2 April 2026}

\begin{document}
\maketitle

\begin{abstract}
We develop the theory of the $\CY_2 \to \CY_3$ fibration: given a $\CY_2$
category $\cC$ (e.g.\ $D^b(\Coh(S))$ for $S$ a K3 surface) and an elliptic
curve $E$, we construct a $\CY_3$ category $\cC_E$ and study the passage from
the quantum vertex chiral group $G_2(\cC)$ to $G_3(\cC_E)$.  The central
mechanism is the \emph{Borcherds multiplicative lift}: the denominator identity
of $G_2(\cC)$ (a weak Jacobi form) lifts to the denominator identity of
$G_3(\cC_E)$ (a Siegel modular form).  We give: (1) the algebraic construction
of the fibered root datum $\cR_3 = \cR_2 \times_E$; (2) the Borcherds lift as
shadow obstruction tower completion; (3) the root multiplicity transfer formula; (4) the
$\Etwo$-structure enhancement from the modular parameter of the elliptic fiber.
The prototypical example is $K3 \times E$, where $\phi_{0,1}$ (the K3 elliptic
genus) lifts to $\Delta_5$ (the Igusa cusp form) via the Borcherds product.
\end{abstract}

\tableofcontents

%% ========================================================================
\section{Introduction and overview}
\label{sec:fib-intro}
%% ========================================================================

A K3 surface $S$ is $\CY_2$; an elliptic curve $E$ is $\CY_1$; the product
$X = S \times E$ is $\CY_3$.  More generally, for an order-$N$ quotient
$X_N = (S \times E)/(\mathbb{Z}/N\mathbb{Z})$, the threefold $X_N$ remains
Calabi--Yau.  The key question:

\medskip
\noindent\textit{How does the quantum vertex chiral group $G_3(X)$ of the $\CY_3$
relate to the quantum vertex chiral group $G_2(S)$ of the $\CY_2$, and what
role does the elliptic curve $E$ play in the passage?}
\medskip

The answer involves three interlocking structures:

\begin{enumerate}[label=(\alph*)]
\item \textbf{Algebraic.}  The generalized root datum $\cR_3(X)$ of the $\CY_3$
  is obtained from the root datum $\cR_2(S)$ of the $\CY_2$ by fibering over
  $E$: one adjoins a new lattice direction (the elliptic fiber class) and
  extends the root system by the data of a Jacobi index.

\item \textbf{Automorphic.}  The denominator identity of $G_2(S)$ is controlled
  by a weak Jacobi form $\phi_{0,m}$ (for K3, $\phi_{0,1}$).  The Borcherds
  multiplicative lift $\mathrm{Borch}(\phi_{0,m})$ produces an automorphic form
  on a larger domain (a Siegel modular form for $\mathrm{Sp}_4(\mathbb{Z})$)
  which is the denominator identity of $G_3(X)$.

\item \textbf{Operadic.}  The $\CY_2$ category already carries an
  $\Etwo$-structure (from $S^2$-framing of cyclic homology).  Fibering over $E$
  does not increase the operad level but introduces a \emph{modular parameter}
  $\tau \in \mathbb{H}$ (the complex structure of $E$) into the $\Etwo$-algebra.
  The resulting object is an $\Etwo$-chiral algebra with an additional
  $\mathrm{SL}_2(\mathbb{Z})$-equivariance from the elliptic fiber.
\end{enumerate}

The passage $G_2 \leadsto G_3$ categorifies the Borcherds lift:
the analytic fact that Jacobi forms lift to Siegel modular forms is the partition-function shadow of an algebraic construction at the level of quantum vertex chiral groups.

%% ========================================================================
\section{The algebraic construction: fibered root data}
\label{sec:fib-algebraic}
%% ========================================================================

\subsection{The $\CY_2$ root datum}
\label{subsec:cy2-root-datum}

Let $\cC$ be a $\CY_2$ category; for concreteness, $\cC = D^b(\Coh(S))$ for
a K3 surface $S$.  The generalized root datum $\cR_2 = \cR_2(\cC)$ consists of:

\begin{definition}[$\CY_2$ root datum]
\label{def:cy2-root-datum}
The $\CY_2$ root datum $\cR_2 = (\Lambda_2, \Delta_2, W_2, \rho_2, \mathrm{mult}_2)$ has:
\begin{enumerate}[label=(\roman*)]
\item \textbf{Lattice.}  $\Lambda_2 = H^*(S, \mathbb{Z})$ with the Mukai
  pairing $\langle \cdot, \cdot \rangle_{\mathrm{Muk}}$.  For K3,
  $\Lambda_2 \simeq U^3 \oplus E_8(-1)^2$, the K3 lattice of signature
  $(4, 20)$.

\item \textbf{Root system.}  $\Delta_2 \subset \Lambda_2$ consists of classes
  $\delta$ with $\langle \delta, \delta \rangle = -2$ (the $(-2)$-curves in
  $S$, i.e.\ the real roots).

\item \textbf{Weyl group.}  $W_2 = W(\Lambda_2)$, generated by reflections in
  the $(-2)$-vectors.

\item \textbf{BPS spectrum.}  For each class $\beta \in \Lambda_2$, the BPS
  count $\mathrm{mult}_2(\beta)$ is determined by the generating function
  \begin{equation}\label{eq:cy2-partition}
    Z_2(\tau, z) = \sum_{\beta} \mathrm{mult}_2(\beta) \, q^{n(\beta)} y^{l(\beta)}
  \end{equation}
  which is a weak Jacobi form of weight $0$ and some index $m$ depending on the
  polarization.  For K3 with the standard polarization, $Z_2 = \phi_{0,1}$.
\end{enumerate}
\end{definition}

\begin{remark}[The K3 elliptic genus as partition function]
\label{rmk:k3-ell-genus}
The weak Jacobi form $\phi_{0,1}(\tau, z)$ of weight $0$ and index $1$ is the
elliptic genus of K3:
\[
  \phi_{0,1}(\tau, z)
  = 2y + 20 + 2y^{-1} + q(20y^2 - 128y + 216 - 128y^{-1} + 20y^{-2}) + \cdots
\]
where $q = e^{2\pi i \tau}$, $y = e^{2\pi i z}$.  Its Fourier coefficients
$f(n, l)$ defined by $\phi_{0,1} = \sum_{n,l} f(n,l) \, q^n y^l$ give the BPS
degeneracies of the K3 sigma model.  This is the ``$\CY_2$ partition function''
that will serve as input to the Borcherds lift.
\end{remark}

\subsection{The elliptic fiber data}
\label{subsec:elliptic-fiber}

Let $(E, \tau_E)$ be an elliptic curve with modular parameter
$\tau_E \in \mathbb{H}$.  The elliptic curve contributes:

\begin{enumerate}[label=(\roman*)]
\item A rank-$2$ lattice $\Lambda_E = H^1(E, \mathbb{Z}) \simeq \mathbb{Z}^2$
  with the symplectic pairing $\langle \alpha, \beta \rangle_E = \alpha \wedge \beta$.

\item The ``fiber direction'' in homology: the class $[E] \in H_2(S \times E)$.

\item The modular parameter $\tau_E$, which promotes the lattice variable $z$
  from the Jacobi form $\phi_{0,m}(\tau, z)$ to a full modular variable.
\end{enumerate}

\subsection{The fibered root datum $\cR_3 = \cR_2 \times_E$}
\label{subsec:fibered-root-datum}

\begin{construction}[Fibered root datum]
\label{constr:fibered-root-datum}
Given a $\CY_2$ root datum
$\cR_2 = (\Lambda_2, \Delta_2, W_2, \rho_2, \mathrm{mult}_2)$
and an elliptic curve $E$, the \emph{fibered root datum}
$\cR_3 = \cR_2 \times_E$ is defined as follows.

\medskip
\noindent\textbf{Step 1: Lattice extension.}  Set
\begin{equation}\label{eq:lattice-extension}
  \Lambda_3 = \Lambda_{\mathrm{hyp}} \oplus \Lambda_E^\vee
\end{equation}
where $\Lambda_{\mathrm{hyp}} \subset \Lambda_2$ is a chosen hyperbolic sublattice
of signature $(2,1)$ and $\Lambda_E^\vee$ contributes one additional direction.
For $K3 \times E$, the relevant sublattice is $\Lambda^{2,1}_{II}$ and the total
lattice is $\Lambda^{3,2}$ of signature $(3,2)$.

Concretely, introduce coordinates $(z_1, z_2, z_3)$ on $\Lambda_3 \otimes \mathbb{C}$
where:
\begin{itemize}
\item $z_1 = \tau$ (the modular parameter, from the ``base'' of the K3 elliptic
  fibration or equivalently from the $\CY_2$ moduli),
\item $z_2 = z$ (the elliptic variable, from the Jacobi form),
\item $z_3 = \sigma$ (the new variable from the elliptic curve $E$).
\end{itemize}
The Siegel upper half-space $\mathbb{H}_2$ is parameterized by the $2 \times 2$
period matrix
\[
  Z = \begin{pmatrix} \tau & z \\ z & \sigma \end{pmatrix}
  = \begin{pmatrix} z_1 & z_2 \\ z_2 & z_3 \end{pmatrix} \in \mathbb{H}_2.
\]

\medskip
\noindent\textbf{Step 2: Root extension.}  The real roots $\Delta_3^{\mathrm{re}}$
are the primitive elements of the fundamental polyhedron $\mathcal{P}$ in the
positive cone of $\Lambda_3$.  For $K3 \times E$, these are the three real simple
roots $\delta_1, \delta_2, \delta_3$ of the type II sublattice $\Lambda^{2,1}_{II}$
with Gram matrix
\[
  (\delta_i, \delta_j) = \begin{pmatrix} 2 & -2 & -2 \\ -2 & 2 & -2 \\ -2 & -2 & 2 \end{pmatrix}.
\]

The imaginary roots $\Delta_3^{\mathrm{im}}$ are indexed by triples $(n, l, m)$
in the positive cone $\{(n, l, m) : (n, l, m) > 0\}$ (with the lexicographic
ordering: $n > 0$, or $n = 0$ and $l > 0$, or $n = l = 0$ and $m > 0$).
Their multiplicities are:
\begin{equation}\label{eq:mult3-from-cy2}
  \mathrm{mult}_3(n, l, m) = f(nm, l)
\end{equation}
where $f(nm, l)$ are the Fourier coefficients of the $\CY_2$ partition function
$\phi_{0,1}$.  This is the \emph{root multiplicity transfer formula}.

\medskip
\noindent\textbf{Step 3: Weyl group extension.}  The Weyl group of $\cR_3$ is
the $2$-reflection group $W_3 = W^{(2)}(\Lambda^{2,1}_{II})$, acting on the
hyperbolic sublattice.  Its structure is
\[
  \mathrm{O}(\Lambda^{2,1})_+ \simeq W^{(2)}(\Lambda^{2,1}_{II}) \rtimes S_3.
\]

\medskip
\noindent\textbf{Step 4: Weyl vector.}  The Weyl vector $\rho_3$ satisfies
$(\rho_3, \delta_i) = -1$ for all real simple roots.  For $K3 \times E$:
\[
  \rho_3 = f_2 - \tfrac{1}{2} f_3 + f_{-2}.
\]
\end{construction}

\begin{remark}[Interpretation: ``Higgs on K3 fibered over $E$'']
\label{rmk:higgs-on-k3}
The construction $\cR_3 = \cR_2 \times_E$ has a direct geometric reading.
The $\CY_2$ datum $\cR_2$ encodes the Higgs branch data: the root system, the $(-2)$-curves, and the BPS spectrum of the K3 sigma model.  Fibering over $E$ promotes each BPS state of charge $\beta \in \Lambda_2$ and spin $l$ to a family parameterized by the winding number $m$ around $E$.  The combined charge $(n, l, m)$ in $\Lambda_3$ records: $n$ = instanton number (base), $l$ = spin/charge (K3 elliptic fiber), $m$ = winding around $E$.

Categorically, this is:
\[
  \cC_E := D^b(\Coh(S \times E)) \simeq \Coh(S) \boxtimes \Coh(E).
\]
The Hochschild homology decomposes as
\[
  \HH_*(\cC_E) \simeq \HH_*(S) \otimes \HH_*(E)
\]
and the $\CY_3$ trace on $\cC_E$ is the product of the $\CY_2$ trace on $S$ and
the $\CY_1$ trace on $E$.  The elliptic fiber contributes the ``winding mode''
quantum number $m$, which, together with the instanton number $n$ from the base,
creates the full two-dimensional lattice indexing the Fourier expansion of a
Siegel modular form.
\end{remark}

\subsection{The fibered root datum in coordinates}
\label{subsec:fibered-coordinates}

To make the construction fully explicit, we spell out how the $\CY_2$ Jacobi
form data embeds into the $\CY_3$ Siegel modular form data.

A weak Jacobi form $\phi_{k,m}(\tau, z)$ of weight $k$ and index $m$ has a
Fourier expansion
\[
  \phi_{k,m}(\tau, z) = \sum_{\substack{n \geq 0 \\ l \in \mathbb{Z} \\ 4nm - l^2 \geq -m^2}} f(n, l) \, q^n y^l,
  \qquad q = e^{2\pi i \tau}, \; y = e^{2\pi i z}.
\]
A Siegel modular form $F(Z)$ on $\mathbb{H}_2$ has the Fourier--Jacobi expansion
\begin{equation}\label{eq:fourier-jacobi}
  F(Z) = F\!\begin{pmatrix} \tau & z \\ z & \sigma \end{pmatrix}
  = \sum_{m \geq 0} \phi_m(\tau, z) \, p^m,
  \qquad p = e^{2\pi i \sigma},
\end{equation}
where each $\phi_m(\tau, z)$ is a Jacobi form of index $m$.

The fibered root datum construction amounts to the following dictionary:

\begin{center}
\begin{tabular}{lll}
\textbf{$\CY_2$ datum} & \textbf{Fibered over $E$} & \textbf{$\CY_3$ datum} \\
\hline
Lattice $\Lambda_2^{\mathrm{hyp}}$ & $\oplus \; \Lambda_E^\vee$ & $\Lambda_3$ \\
Jacobi form $\phi_{0,m}(\tau, z)$ & Fourier--Jacobi coeff.\ of & Siegel form $F(Z)$ \\
Fourier coeff.\ $f(n, l)$ & $\times$ winding $m$ & mult$_3(n,l,m) = f(nm,l)$ \\
Modular group $\mathrm{SL}_2(\mathbb{Z})$ & $\times$ fiber $\mathrm{SL}_2(\mathbb{Z})$ & Siegel modular group $\mathrm{Sp}_4(\mathbb{Z})$ \\
$\tau \in \mathbb{H}$ & $\times \; \sigma \in \mathbb{H}$ & $Z \in \mathbb{H}_2$
\end{tabular}
\end{center}

%% ========================================================================
\section{The Borcherds lift as shadow obstruction tower completion}
\label{sec:fib-borcherds}
%% ========================================================================

\subsection{The Borcherds multiplicative lift}
\label{subsec:borcherds-lift}

We recall the Borcherds multiplicative lift and interpret it via the shadow obstruction tower of Volume~I.

\begin{theorem}[Borcherds]
\label{thm:borcherds-lift}
Let $\phi(\tau, z) = \sum_{n, l} f(n, l) \, q^n y^l$ be a weak Jacobi form of
weight $0$ and index $1$ for $\mathrm{SL}_2(\mathbb{Z})$ with
$f(n, l) \in \mathbb{Z}$ and $f(n, l)$ depending only on $4n - l^2$ and $l \bmod 2$.
The Borcherds lift $\mathrm{Borch}(\phi)$ is the meromorphic Siegel modular form
of weight $f(0,0)/2$ defined by
\begin{equation}\label{eq:borcherds-product}
  \mathrm{Borch}(\phi)(Z)
  = q^A y^B p^C \prod_{\substack{(n,l,m) \in \mathbb{Z}^3 \\ (n,l,m) > 0}}
  \bigl(1 - q^n y^l p^m\bigr)^{f(nm, l)}
\end{equation}
where $q = e^{2\pi i \tau}$, $y = e^{2\pi i z}$, $p = e^{2\pi i \sigma}$,
$Z = \bigl(\begin{smallmatrix} \tau & z \\ z & \sigma \end{smallmatrix}\bigr)$,
and $A, B, C$ are determined by the Weyl vector $\rho_3$.

The product converges on an appropriate domain and extends to an automorphic form
for a subgroup of $\mathrm{Sp}_4(\mathbb{Z})$ (the full group if $\phi$ has
level $1$).
\end{theorem}

\begin{example}[The Igusa cusp form $\Delta_5$]
\label{ex:igusa-from-borcherds}
The K3 elliptic genus $\phi_{0,1}$ has $f(0, 0) = 20$, so the Borcherds lift
has weight $20/2 = 10$.  After the rescaling $Z \mapsto 2Z$ and the prefactor
$1/64$, one obtains the Igusa cusp form of weight $5$:
\begin{align}\label{eq:delta5-product}
  \frac{1}{64} \Delta_5(2Z)
  &= e^{-2\pi i \langle \rho_3, z \rangle}
  \prod_{\substack{(n,l,m) > 0}} \bigl(1 - e^{2\pi i(n\tau + lz + m\sigma)}\bigr)^{f(nm, l)} \\
  &= \mathrm{Borch}(\phi_{0,1})(2Z) / 64. \notag
\end{align}
This is the denominator identity of the BKM superalgebra
$\mathfrak{g}_{\Delta_5}$.
\end{example}

\subsection{Shadow obstruction tower interpretation}
\label{subsec:shadow-tower-interp}

The Borcherds lift admits the following interpretation in terms of the shadow
tower of Volume~I.

\begin{construction}[Shadow obstruction tower from $\CY_2$ to $\CY_3$]
\label{constr:shadow-tower-cy2-cy3}
Let $A_2 = \Phi(\cC)$ be the quantum chiral algebra of the $\CY_2$ category
$\cC$, with shadow obstruction tower $\Theta_{A_2} = \sum_{r \geq 2} \Theta^{(r)}_{A_2}$.
The $\CY_3$ shadow obstruction tower $\Theta_{A_3}$ of $A_3 = \Phi(\cC_E)$ is the
\emph{shadow obstruction tower completion} of $\Theta_{A_2}$ along the elliptic fiber:

\begin{enumerate}[label=(\roman*)]
\item \textbf{Arity 2 (the bilinear shadow, $\kappa_{\mathrm{ch}}$).}  The leading term
  $\Theta^{(2)}_{A_3}$ is the modular characteristic $\kappa_{\mathrm{ch}}(A_3)$, which equals
  the weight of $\Delta_5$:
  \[
    \kappa_{\mathrm{ch}}(A_3) = 5.
  \]
  This arises from $\kappa_{\mathrm{ch}}(A_2)$ (the $\CY_2$ modular characteristic) via
  the relation
  \[
    \kappa_{\mathrm{ch}}(A_3) = \frac{f(0,0)}{2} \cdot \frac{1}{2} = \frac{20}{4} = 5
  \]
  where $f(0,0) = 20$ is the constant term of $\phi_{0,1}$ and the factor
  $1/2$ accounts for the rescaling $Z \mapsto 2Z$.

\item \textbf{Arity 3 (the cubic shadow, $C$).}  The first correction to the
  Kac--Moody approximation adds the simplest imaginary roots.  In the Borcherds
  product, these are the factors with $nm = 0$ (i.e., either $n = 0$ or $m = 0$):
  \[
    \Theta^{(3)}_{A_3} \;\longleftrightarrow\; \prod_{\substack{(n,l,m) > 0 \\ nm = 0}}
    (1 - q^n y^l p^m)^{f(0, l)}.
  \]
  Since $f(0, l)$ for $\phi_{0,1}$ gives $f(0, 0) = 20$, $f(0, \pm 1) = 2$,
  this is the contribution of the purely ``fiber'' or purely ``base'' imaginary
  roots, those that do not mix the $\tau$ and $\sigma$ directions.

\item \textbf{Arity $r$ (the higher shadows).}  The arity-$r$ shadow
  $\Theta^{(r)}_{A_3}$ adds imaginary roots with $nm \leq r - 2$ (charge mixing
  up to level $r - 2$).  In the limit $r \to \infty$, the full Borcherds product
  is recovered:
  \[
    \Theta_{A_3} = \lim_{r \to \infty} \Theta^{(\leq r)}_{A_3}
    \quad \longleftrightarrow \quad
    \text{full infinite product in \eqref{eq:borcherds-product}}.
  \]
\end{enumerate}

The convergence of the Borcherds product (in a suitable domain of $\mathbb{H}_2$)
is thus the analytic manifestation of the algebraic completeness of the shadow
tower.
\end{construction}

\begin{remark}[The denominator identity as bar complex Euler product]
\label{rmk:denom-bar-euler}
In Volume~I, the denominator identity of a BKM superalgebra is the Euler characteristic of the bar complex $B(A)$.  The Borcherds product formula
\[
  \Phi(z) = e^{-2\pi i \langle \rho, z \rangle}
  \prod_{\alpha \in \Delta_+} (1 - e^{-2\pi i \langle \alpha, z \rangle})^{\mathrm{mult}\, \alpha}
\]
is literally the ``$q$-graded Euler product'' of the bar complex: each factor
$(1 - e^{-2\pi i \langle \alpha, z \rangle})$ corresponds to a generator of
$B(A)$ in root degree $\alpha$, raised to the power $\mathrm{mult}\, \alpha$
(the number of independent generators in that degree).

The Borcherds lift says: the bar-complex Euler product of the $\CY_3$ algebra $A_3$ is determined by the Fourier coefficients of the $\CY_2$ Euler product.  This is the fibered bar complex structure.
\end{remark}

\subsection{The Borcherds lift in four steps}
\label{subsec:borcherds-four-steps}

We now present the Borcherds lift as a four-step algebraic procedure, making its shadow-tower nature manifest.

\begin{construction}[Borcherds lift as algebraic procedure]
\label{constr:borcherds-algebraic}
\

\medskip
\noindent\textbf{Step 1: Extract the $\CY_2$ input.}
Start with the $\CY_2$ partition function $\phi = \phi_{0,1}(\tau, z)$, a weak
Jacobi form.  Its Fourier coefficients $\{f(n, l)\}_{n \geq 0, l \in \mathbb{Z}}$
are the $\CY_2$ root multiplicities.

\medskip
\noindent\textbf{Step 2: Construct the $\CY_3$ root multiplicities.}
For each triple $(n, l, m)$ in the positive cone, define
\[
  \mathrm{mult}_3(n, l, m) := f(nm, l).
\]
This is the root multiplicity transfer formula.  The ``mixing'' of $n$ and $m$
via the product $nm$ is the essential feature: the $\CY_2$ Fourier coefficient
at discriminant $D = 4nm - l^2$ controls the $\CY_3$ root multiplicity at
charge $(n, l, m)$.

\medskip
\noindent\textbf{Step 3: Form the automorphic product.}
Construct the infinite product
\[
  \Phi_3(Z) = e^{-2\pi i \langle \rho_3, z \rangle}
  \prod_{(n,l,m) > 0} (1 - q^n y^l p^m)^{\mathrm{mult}_3(n,l,m)}.
\]
The Borcherds theorem guarantees this converges (in a suitable domain) and
defines an automorphic form for $\mathrm{Sp}_4(\mathbb{Z})$.

\medskip
\noindent\textbf{Step 4: Verify the denominator identity.}
The BKM superalgebra $\mathfrak{g}_3$ with root multiplicities $\mathrm{mult}_3$
has the Weyl--Kac--Borcherds denominator identity
\[
  \sum_{w \in W_3} \det(w) \, e^{-\pi i \langle w(\rho_3), z \rangle}
  = e^{-2\pi i \langle \rho_3, z \rangle}
  \prod_{\alpha \in \Delta_{3,+}} (1 - e^{-2\pi i \langle \alpha, z \rangle})^{\mathrm{mult}_3(\alpha)}.
\]
The left side is a sum over the Weyl group $W_3$; the right side is the
Borcherds product.  Both sides equal $\frac{1}{64} \Delta_5(2Z)$.
\end{construction}

%% ========================================================================
\section{Root multiplicity transfer}
\label{sec:fib-root-mult}
%% ========================================================================

\subsection{The transfer formula}
\label{subsec:transfer-formula}

The root multiplicity transfer is the heart of the $\CY_2 \to \CY_3$ passage.

\begin{proposition}[Root multiplicity transfer]
\label{prop:root-mult-transfer}
Let $\cC$ be a $\CY_2$ category with partition function
$\phi_{0,m}(\tau, z) = \sum f(n, l) \, q^n y^l$, and let $\cC_E$ be the
fibered $\CY_3$ category.  The root multiplicities of the $\CY_3$ algebra
$\mathfrak{g}_3 = \mathfrak{g}(\cC_E)$ are:
\begin{equation}\label{eq:root-mult-transfer}
  \mathrm{mult}_3(\alpha) = f\bigl(\langle \alpha, \alpha \rangle_{\tau\sigma}, \;
  \langle \alpha, \alpha \rangle_z\bigr)
\end{equation}
where, in the coordinates $\alpha = (n, l, m)$:
\begin{itemize}
\item $\langle \alpha, \alpha \rangle_{\tau\sigma} = nm$ (the product of the
  $\tau$- and $\sigma$-components),
\item $\langle \alpha, \alpha \rangle_z = l$ (the $z$-component).
\end{itemize}
More explicitly: $\mathrm{mult}_3(n, l, m) = f(nm, l)$.
\end{proposition}

\begin{proof}[Proof sketch]
This follows from the Borcherds product formula.  The product side of the
denominator identity of $\mathfrak{g}_3$ is
\[
  \prod_{(n,l,m) > 0} (1 - q^n y^l p^m)^{c(n,l,m)}
\]
where $c(n, l, m) = \mathrm{mult}_3(n, l, m)$.  Comparing with the Borcherds
lift of $\phi$, which has exponents $f(nm, l)$, gives
$c(n, l, m) = f(nm, l)$.
\end{proof}

\subsection{Consequences of the transfer formula}
\label{subsec:transfer-consequences}

\begin{corollary}[Discriminant dependence]
\label{cor:discriminant-dependence}
The root multiplicity $\mathrm{mult}_3(n, l, m)$ depends on $n$ and $m$ only
through the product $nm$.  Equivalently, $\mathrm{mult}_3$ depends on the
``discriminant'' $D = 4nm - l^2$ and the parity of $l$.  This is because
$f(nm, l)$, for $\phi_{0,1}$, depends on $nm$ and $l$ only through
$4(nm) - l^2$ and $l \bmod 2$.
\end{corollary}

\begin{proof}
This is a standard property of Jacobi forms of index $1$: the Fourier
coefficient $f(n, l)$ depends only on $4n - l^2$ and $l \bmod 2$.  The
substitution $n \mapsto nm$ gives the result.
\end{proof}

\begin{corollary}[Real root multiplicities]
\label{cor:real-root-mult}
The real roots of $\mathfrak{g}_3$ (those $\alpha$ with $(\alpha, \alpha) = 2$)
all have multiplicity $1$.  This corresponds to $f(nm, l)$ at $4nm - l^2 = -1$,
which gives $f = 1$ for $\phi_{0,1}$.
\end{corollary}

\begin{corollary}[Lightlike root multiplicities]
\label{cor:lightlike-mult}
The lightlike (isotropic) roots with $(\alpha, \alpha) = 0$ have multiplicity
$f(0, 0) = 20$ when $l$ is even, and $f(0, 1) = 2$ when $l$ is odd.  These are
the ``threshold'' imaginary roots.  For K3, the multiplicity 20 corresponds to
the 20 directions in $H^{1,1}(S)$ transverse to the K\"ahler class.
\end{corollary}

\subsection{The multiplicity table}
\label{subsec:mult-table}

For the K3 case ($\phi_{0,1}$), the first few root multiplicities are:

\begin{center}
\begin{tabular}{c|ccccc}
$nm \;\backslash\; l$ & $0$ & $\pm 1$ & $\pm 2$ & $\pm 3$ & $\pm 4$ \\
\hline
$0$ & $20$ & $2$ & $0$ & $0$ & $0$ \\
$1$ & $216$ & $-128$ & $20$ & $2$ & $0$ \\
$2$ & $1616$ & $-1026$ & $216$ & $2$ & $0$ \\
$3$ & $8032$ & $-5504$ & $1616$ & $-128$ & $0$ \\
$4$ & $32120$ & $-23550$ & $8032$ & $-1026$ & $20$
\end{tabular}
\end{center}

Negative multiplicities correspond to fermionic (odd) generators:
$\mathrm{mult}_3 < 0$ means $|\mathrm{mult}_3|$ fermionic generators in the
corresponding root space.  This is the super-structure of the BKM superalgebra.

\begin{remark}[Physical interpretation]
\label{rmk:physical-interp}
The multiplicity table encodes the BPS spectrum of the type IIA string on
$K3 \times E$:
\begin{itemize}
\item $nm$ = D0-D2 charge (instanton number times D2 wrapping),
\item $l$ = D2 charge along the elliptic fiber of K3,
\item $f(nm, l) > 0$: bosonic BPS states,
\item $f(nm, l) < 0$: fermionic BPS states.
\end{itemize}
The root multiplicity transfer formula says that the $\CY_3$ BPS spectrum is
entirely determined by the $\CY_2$ BPS spectrum (the K3 elliptic genus),
``fibered'' over the additional winding modes from $E$.
\end{remark}

%% ========================================================================
\section{The $\Etwo$-structure and the modular parameter}
\label{sec:fib-e2}
%% ========================================================================

\subsection{$\Etwo$ from $\CY_2$}
\label{subsec:e2-from-cy2}

A $\CY_2$ category $\cC$ carries a natural $\Etwo$-structure.  The origin is
the $S^2$-framing of the cyclic homology:

\begin{proposition}[$\CY_2$ gives $\Etwo$]
\label{prop:cy2-e2}
Let $\cC$ be a smooth proper $\CY_2$ category.  The Hochschild cohomology
$\HH^*(\cC)$ carries an $\Etwo$-algebra structure (equivalently, a
Gerstenhaber algebra structure via Kontsevich formality).  The associated
quantum chiral algebra $A_2 = \Phi(\cC)$ is an $\Etwo$-chiral algebra, and
$\Rep^{\Etwo}(A_2)$ is a braided monoidal category.
\end{proposition}

This $\Etwo$-structure is the ``seed'' for the quantum group structure.  At the
level of the root datum $\cR_2$, it manifests as:
\begin{itemize}
\item The modular transformation $\tau \mapsto -1/\tau$ is the braiding
  (the $S$-matrix of the modular tensor category),
\item The $T$-transformation $\tau \mapsto \tau + 1$ is the twist.
\end{itemize}
The modular group $\mathrm{SL}_2(\mathbb{Z})$ acting on $\tau \in \mathbb{H}$
is the ``$\CY_2$-level'' manifestation of the $\Etwo$-braiding.

\subsection{Enhancement by the elliptic fiber}
\label{subsec:e2-enhancement}

Fibering over $E$ introduces a second modular parameter $\sigma$.  The full
modular group upgrades from $\mathrm{SL}_2(\mathbb{Z})$ to
$\mathrm{Sp}_4(\mathbb{Z})$, acting on $Z \in \mathbb{H}_2$.

\begin{proposition}[$\Etwo$ with modular parameter]
\label{prop:e2-modular}
The $\CY_3$ quantum chiral algebra $A_3 = \Phi(\cC_E)$ is an $\Etwo$-chiral
algebra whose $\Etwo$-structure is enhanced by a modular parameter: the braided
monoidal category $\Rep^{\Etwo}(A_3)$ carries a compatible
$\mathrm{Sp}_4(\mathbb{Z})$-action.

More precisely, the data organizing the passage from $\Etwo$ over
$\mathrm{SL}_2(\mathbb{Z})$ to $\Etwo$ over $\mathrm{Sp}_4(\mathbb{Z})$ is:
\begin{enumerate}[label=(\roman*)]
\item The Jacobi group $\mathrm{SL}_2(\mathbb{Z}) \ltimes \mathbb{Z}^2$ acts on
  $(\tau, z)$ and controls the Jacobi form $\phi_{0,1}(\tau, z)$.  This is the
  ``$\CY_2$ modular structure.''

\item The Siegel modular group $\mathrm{Sp}_4(\mathbb{Z})$ acts on
  $Z = \bigl(\begin{smallmatrix} \tau & z \\ z & \sigma \end{smallmatrix}\bigr)$
  and controls the Siegel form $\Delta_5(Z)$.  This is the ``$\CY_3$ modular
  structure.''

\item The Jacobi group embeds into $\mathrm{Sp}_4(\mathbb{Z})$ via the
  ``Fourier--Jacobi parabolic,'' the stabilizer of the cusp
  $\sigma \to i\infty$.  The Fourier--Jacobi expansion
  \[
    \Delta_5(Z) = \sum_{m \geq 1} \phi_m(\tau, z) \, p^m
  \]
  is the decomposition of the $\mathrm{Sp}_4(\mathbb{Z})$-representation
  into Jacobi group representations.
\end{enumerate}
\end{proposition}

\begin{remark}[The modular parameter is \emph{not} an extra $\Eone$]
\label{rmk:not-extra-e1}
A natural expectation is that the passage $\CY_2 \to \CY_3$ upgrades $\Etwo$ to
$E_3$ (by Dunn: $E_3 \simeq \Etwo \otimes \Eone$).  This is \emph{not} what
happens.  The $\CY_3$ manifold $S \times E$ has trivial canonical bundle of
dimension $3$, giving an $S^3$-framing, and $E_3$ does appear at the topological
level.  But the \emph{chiral} (holomorphic) structure only sees $\Etwo$: by
the Dunn decomposition $E_3 \to \Etwo$, one $\Eone$-factor is ``used up'' by
the holomorphic structure of $E$.  What remains is an $\Etwo$-chiral algebra
with an additional $\mathrm{SL}_2(\mathbb{Z})$-equivariance parametrized by
$\sigma$.

Alternatively: the ``third direction'' from $E$ does not add an independent
operadic degree of freedom, but rather introduces a \emph{spectral} parameter
$\sigma$ into the existing $\Etwo$-structure.  The $\Etwo$-braiding $R(z)$ of
the $\CY_2$ algebra becomes a family $R(z; \sigma)$ of braidings parametrized
by $\sigma \in \mathbb{H}$.  This family is controlled by the Fourier--Jacobi
expansion.
\end{remark}

\subsection{The Sp$_4$ action as $\Etwo$ braiding}
\label{subsec:sp4-e2}

The $\mathrm{Sp}_4(\mathbb{Z})$-action on $\mathbb{H}_2$ decomposes into
generators:

\begin{align}
  S &\colon Z \longmapsto -Z^{-1}
  &&\text{(braiding / $S$-matrix)}, \label{eq:sp4-S} \\
  T &\colon Z \longmapsto Z + \bigl(\begin{smallmatrix} 1 & 0 \\ 0 & 0 \end{smallmatrix}\bigr)
  &&\text{(twist in $\tau$-direction)}, \label{eq:sp4-T} \\
  U &\colon Z \longmapsto Z + \bigl(\begin{smallmatrix} 0 & 0 \\ 0 & 1 \end{smallmatrix}\bigr)
  &&\text{(twist in $\sigma$-direction)}, \label{eq:sp4-U} \\
  V &\colon Z \longmapsto Z + \bigl(\begin{smallmatrix} 0 & 1 \\ 1 & 0 \end{smallmatrix}\bigr)
  &&\text{(mixed twist)}. \label{eq:sp4-V}
\end{align}

The generators $T$ and $U$ are the two ``$\Eone$-twists'' (one from the $\CY_2$
base, one from the elliptic fiber).  The generator $S$ is the $\Etwo$-braiding.
The generator $V$ is the \emph{mixing} of the two $\Eone$-structures, which is
the genuinely new datum of the fibered construction.

\begin{proposition}[Jacobi form as $\Etwo$-character]
\label{prop:jacobi-e2-char}
The weak Jacobi form $\phi_{0,1}(\tau, z)$ is the character (graded trace) of
the $\CY_2$ quantum chiral algebra $A_2$:
\[
  \phi_{0,1}(\tau, z) = \Tr_{A_2}\bigl(q^{L_0} y^{J_0}\bigr)
\]
where $L_0$ is the Virasoro zero-mode (Hamiltonian) and $J_0$ is the
$\mathrm{U}(1)$-charge (the Jacobi index variable).

The Siegel modular form $\Delta_5(Z)$ is the denominator of the $\CY_3$
algebra $A_3$:
\[
  \Delta_5(Z) \;\sim\; \det\nolimits_{\frakg_3}\bigl(1 - e^{2\pi i Z}\bigr)
  := \prod_{\alpha > 0} (1 - e^{2\pi i \langle \alpha, Z \rangle})^{\mathrm{mult}_3(\alpha)}.
\]
The Borcherds lift is the passage from the character of $A_2$ to the
denominator of $A_3$.
\end{proposition}

%% ========================================================================
\section{Generalization and structural theorems}
\label{sec:fib-general}
%% ========================================================================

\subsection{General $\CY_2$ categories}
\label{subsec:general-cy2}

The construction of Section~\ref{sec:fib-algebraic} generalizes beyond
$D^b(\Coh(S))$ for K3 surfaces.

\begin{conjecture}[General $\CY_2 \to \CY_3$ fibration]
\label{conj:general-fibration}
Let $\cC$ be any smooth proper $\CY_2$ category with partition function
$\phi_\cC(\tau, z)$ (a weak Jacobi form of weight $0$ and some index $m$).
Then:
\begin{enumerate}[label=(\roman*)]
\item The fibered category $\cC_E = \cC \boxtimes D^b(\Coh(E))$ is a $\CY_3$
  category.

\item The root datum $\cR_3(\cC_E)$ is the fibered root datum
  $\cR_2(\cC) \times_E$ of Construction~\ref{constr:fibered-root-datum}.

\item The denominator identity of the quantum vertex chiral group
  $G_3(\cC_E)$ is the Borcherds lift $\mathrm{Borch}(\phi_\cC)$.

\item The root multiplicities are given by
  $\mathrm{mult}_3(n, l, m) = f_\cC(nm, l)$.

\item The $\Etwo$-braiding of $G_3(\cC_E)$ is the
  $\mathrm{Sp}_4(\mathbb{Z})$-equivariant enhancement of the
  $\mathrm{SL}_2(\mathbb{Z})$-braiding of $G_2(\cC)$.
\end{enumerate}
\end{conjecture}

\begin{remark}
\label{rmk:index-m-general}
For a general $\CY_2$ category, the partition function may be a Jacobi form of
index $m > 1$.  The Borcherds lift of a Jacobi form of index $m$ produces a
Siegel modular form for the paramodular group $\Gamma_m \subset \mathrm{Sp}_4(\mathbb{Q})$
rather than the full $\mathrm{Sp}_4(\mathbb{Z})$.  This reflects the fact that the
elliptic fibration structure of $\cC_E$ may have non-trivial monodromy of
order related to $m$.
\end{remark}

\subsection{The fibration as functor}
\label{subsec:fibration-functor}

\begin{construction}[The fibration functor $(-) \times_E$]
\label{constr:fibration-functor}
The fibration over $E$ defines a functor
\[
  (-) \times_E \;\colon\; \CY_2\text{-}\mathrm{QChirAlg}
  \longrightarrow \CY_3\text{-}\mathrm{QChirAlg}
\]
from the category of $\CY_2$ quantum chiral algebras to $\CY_3$ quantum chiral
algebras.  On objects, it sends $G_2(\cC)$ to $G_3(\cC_E)$.  On morphisms (i.e.,
on homomorphisms of quantum vertex chiral groups), it acts by:
\begin{enumerate}[label=(\roman*)]
\item Extending the root lattice by the fiber direction,
\item Lifting the partition function via Borcherds,
\item Enhancing the modular group from $\mathrm{SL}_2$ to $\mathrm{Sp}_4$.
\end{enumerate}

At the level of shadow obstruction towers:
\[
  \Theta_{G_3} = \mathrm{Borch}\bigl(\exp(\Theta_{G_2})\bigr)
\]
where $\exp(\Theta_{G_2})$ denotes the ``exponentiated shadow obstruction tower'' (the
passage from the Lie-algebraic MC element to the group-like denominator
product) and $\mathrm{Borch}$ is the Borcherds multiplicative lift.
\end{construction}

\subsection{Compatibility with the main theorems}
\label{subsec:compatibility-main}

The fibration functor is compatible with the four target theorems CY-A through
CY-D:

\begin{center}
\begin{tabular}{lll}
\textbf{Theorem} & \textbf{$\CY_2$ version} & \textbf{$\CY_3 = \CY_2 \times_E$ version} \\
\hline
CY-A & $\Phi(\cC) = G_2(\cC)$ & $\Phi(\cC_E) = G_3(\cC_E) = G_2 \times_E$ \\
CY-B & Bar-cobar for $G_2$ & Fibered bar-cobar for $G_3$ \\
CY-C & $\Rep^{\Etwo}(G_2)$ braided & $\Rep^{\Etwo}(G_3)$ braided, $\mathrm{Sp}_4$-equiv.\ \\
CY-D & $\kappa_{\mathrm{ch}}(G_2) = \chi^{\CY}(\cC)$ & $\kappa_{\mathrm{ch}}(G_3) = \mathrm{wt}(\mathrm{Borch}(\phi_\cC))$
\end{tabular}
\end{center}

%% ========================================================================
\section{The eight diagonal divisor forms}
\label{sec:fib-eight}
%% ========================================================================

The $K3 \times E$ construction admits $\mathbb{Z}/N\mathbb{Z}$ quotients.  The
eight diagonal divisor (dd) Siegel modular forms arise as Borcherds lifts of
the eight twisted-twined K3 elliptic genera.

\begin{conjecture}[Eight fibrations]
\label{conj:eight-fibrations}
For each $N \in \{1, 2, 3, 4, 5, 6, 7, 8\}$ (or more precisely, for each
conjugacy class in $M_{24}$ yielding a distinct dd-form), the $\CY_3$
$X_N = (S \times E)/(\mathbb{Z}/N\mathbb{Z})$ satisfies:
\begin{enumerate}[label=(\roman*)]
\item The $\CY_2$ input is the $g_N$-twisted K3 elliptic genus
  $\phi_{g_N}(\tau, z)$.

\item The Borcherds lift $\mathrm{Borch}(\phi_{g_N})$ is the dd-modular form
  $\Delta_{5,N}$.

\item The quantum vertex chiral group $G_3(X_N)$ has denominator identity
  $\Delta_{5,N}$.

\item The root multiplicities of $G_3(X_N)$ are the Fourier coefficients of
  $\phi_{g_N}$.
\end{enumerate}
The eight quantum vertex chiral groups $\{G_3(X_N)\}$ form a family related by
Galois conjugation on the $N$-torsion of $E$.
\end{conjecture}

%% ========================================================================
\section{Iterated fibration and higher CY}
\label{sec:fib-iterated}
%% ========================================================================

\begin{remark}[Iterated fibration]
\label{rmk:iterated}
The construction naturally iterates.  Starting from a $\CY_1$ category (i.e., an
elliptic curve $E_0$), fibering over $E_1$ gives a $\CY_2$ category (an abelian
surface $E_0 \times E_1$), and fibering again over $E_2$ gives a $\CY_3$
category ($E_0 \times E_1 \times E_2$, an abelian threefold).

At each stage:
\begin{itemize}
\item $\CY_1$: the partition function is a modular form (a function on $\mathbb{H}$),
\item $\CY_2$: the partition function is a Jacobi form (a function on $\mathbb{H} \times \mathbb{C}$),
\item $\CY_3$: the partition function is a Siegel modular form (a function on $\mathbb{H}_2$).
\end{itemize}
The pattern suggests that $\CY_d$ partition functions should be Siegel modular
forms of degree $\lfloor d/2 \rfloor$, and the $\CY_d \to \CY_{d+1}$ fibration
should be a generalized Borcherds lift from degree-$\lfloor d/2 \rfloor$ to
degree-$\lfloor (d+1)/2 \rfloor$ forms.  For $d = 3 \to 4$, one expects
Siegel modular forms of degree $2$ lifting to automorphic forms on an
orthogonal group of higher rank.
\end{remark}

\begin{conjecture}[Borcherds tower]
\label{conj:borcherds-tower}
There exists a sequence of ``Borcherds lift'' functors
\[
  \CY_1\text{-}\mathrm{QChirAlg}
  \xrightarrow{\;(-) \times_{E_1}\;}
  \CY_2\text{-}\mathrm{QChirAlg}
  \xrightarrow{\;(-) \times_{E_2}\;}
  \CY_3\text{-}\mathrm{QChirAlg}
  \xrightarrow{\;(-) \times_{E_3}\;}
  \cdots
\]
such that the passage from level $d$ to level $d+1$ is controlled by a
generalized Borcherds lift.  At each level, the denominator identity of the
quantum vertex chiral group is an automorphic form of the appropriate type.
\end{conjecture}

%% ========================================================================
\section{Open questions}
\label{sec:fib-open}
%% ========================================================================

\begin{enumerate}[label=\arabic*.]
\item \textbf{Non-product CY$_3$s.}  The fibration $\cC_E = \cC \boxtimes
  \Coh(E)$ is a product.  What about $\CY_3$ manifolds that are non-trivially
  fibered over a curve?  For example, an elliptic fibration $X \to B$ where
  $B$ is a rational or higher-genus surface.  The Borcherds lift should
  generalize, but the input may no longer be a standard Jacobi form.

\item \textbf{The inverse problem.}  Given a $\CY_3$ quantum vertex chiral group
  $G_3$, when does it ``defibrate,'' i.e., when is there a $\CY_2$ category
  $\cC$ and an elliptic curve $E$ such that $G_3 = G_2(\cC) \times_E$?  This
  is the question of when a Siegel modular form is a Borcherds lift.  The answer
  involves the ``singular theta correspondence'' of Borcherds--Harvey--Moore.

\item \textbf{Wall-crossing.}  The root multiplicities $f(nm, l)$ are the
  ``indexed'' BPS counts.  How does wall-crossing on the $\CY_3$ side
  (Kontsevich--Soibelman) relate to wall-crossing on the $\CY_2$ side?  The
  fibration should intertwine the two wall-crossing structures.

\item \textbf{Koszul duality of the fibration.}  Is the Koszul dual of
  $G_3(\cC_E)$ related to $G_2(\cC^!)_E$, where $\cC^!$ is the Koszul dual
  $\CY_2$ category?

\item \textbf{Categorical Borcherds lift.}  Can the Borcherds multiplicative
  lift be categorified to a functor between appropriate derived categories?
  I.e., is there a categorical lift
  $\mathrm{Borch} \colon D^b(\Rep^{\Etwo}(G_2)) \to D^b(\Rep^{\Etwo}(G_3))$?

\item \textbf{Connection to Mathieu moonshine.}  The eight twisted K3 elliptic
  genera and the role of $M_{24}$ in the K3 elliptic genus suggest a connection
  between the fibration theory and Mathieu moonshine.  Can the quantum vertex
  chiral group framework illuminate why $M_{24}$ appears?
\end{enumerate}

%% ========================================================================
\section{Summary: the fibration dictionary}
\label{sec:fib-summary}
%% ========================================================================

\begin{center}
\renewcommand{\arraystretch}{1.3}
\begin{tabular}{lll}
\textbf{$\CY_2$ side} & \textbf{Fibration} & \textbf{$\CY_3$ side} \\
\hline
Category $\cC$ & $\boxtimes \Coh(E)$ & Category $\cC_E$ \\
Root datum $\cR_2$ & $\times_E$ & Root datum $\cR_3$ \\
Lattice $\Lambda_2^{\mathrm{hyp}}$ & $\oplus \, \Lambda_E^\vee$ & Lattice $\Lambda_3$ \\
Jacobi form $\phi_{0,m}$ & Borcherds lift & Siegel form $\Delta_k$ \\
$f(n, l)$ & $(n,l) \mapsto (nm, l)$ & $\mathrm{mult}_3(n,l,m) = f(nm, l)$ \\
$\mathrm{SL}_2(\mathbb{Z})$ & Fourier--Jacobi embed.\ & $\mathrm{Sp}_4(\mathbb{Z})$ \\
$\tau \in \mathbb{H}$ & $\times \, \sigma$ & $Z \in \mathbb{H}_2$ \\
$\Etwo$-braiding $R(z)$ & param.\ by $\sigma$ & $\Etwo$-braiding $R(z; \sigma)$ \\
Shadow obstruction tower $\Theta_{G_2}$ & $\mathrm{Borch}(\exp(-))$ & Shadow obstruction tower $\Theta_{G_3}$ \\
$\kappa_{\mathrm{ch}}(G_2) = \chi^{\CY}(\cC)$ & $\times \frac{1}{2} f(0,0)$ & $\kappa_{\mathrm{ch}}(G_3) = \mathrm{wt}(\Delta_k)$ \\
Bar complex $B(A_2)$ & Fibered bar complex & Bar complex $B(A_3)$ \\
$\Eone$-sector (CoHA) & $\times_E$ & $\Eone$-sector (CoHA of $\cC_E$) \\
\end{tabular}
\end{center}

%% ========================================================================
\section{The relative construction: K3 fibrations and shadow depth upgrade}
\label{sec:relative-construction}
%% ========================================================================

The swarm computation (81~tests) develops the \emph{relative}
version of the $\CY_2 \to \CY_3$ fibration: starting from
a K3-fibered CY3, construct the chiral algebra and shadow
tower by relativizing over the base.

\subsection{K3 fibration and relative chiral algebra}

Let $\pi \colon X \to B$ be an elliptically fibered K3 surface
$S$ fibered over a curve $B$, giving a CY3 $X$ with
$\pi \colon X \to B$ a K3 fibration.  The relative construction:

\begin{enumerate}[label=(\roman*)]
  \item The \textbf{fiber chiral algebra} $A_S$ is the $\CY_2$
    algebra of the K3 surface, with $\kappa_{\mathrm{ch}}(A_S) = 24$ (the
    Euler characteristic of K3).

  \item The \textbf{relative chiral algebra} $A_{X/B}$ is
    defined fiberwise: over each point $b \in B$, the fiber
    $X_b = S$ contributes $A_S$, and the variation over $B$
    gives a family of chiral algebras parametrized by moduli.

  \item The \textbf{global chiral algebra} $A_X$ of the CY3
    has $\kappa_{\mathrm{ch}}(A_X) = 5$ (the weight of the Igusa cusp form
    $\Delta_5$).  This is \emph{not} $\kappa_{\mathrm{ch}}(A_S) = 24$:
    the passage from fiber to total space changes $\kappa_{\mathrm{ch}}$
    drastically.
\end{enumerate}

\subsection{Shadow tower of the fibration}

The shadow tower of $A_X$ receives contributions from two
sources: the fiber shadow (from $A_S$) and the base shadow
(from the variation of $A_S$ over $B$).

\begin{itemize}
  \item \textbf{Fiber shadow}: the K3 algebra $A_S$ is
    class $G$ (Gaussian, $r_{\max} = 2$) when viewed as a
    $\CY_2$ algebra.  Its shadow tower terminates at arity~2:
    the only invariant is $\kappa_{\mathrm{ch}} = 24$.

  \item \textbf{Global shadow}: the CY3 algebra $A_X$ is
    class $M$ (mixed, $r_{\max} = \infty$).  The infinite
    tower of imaginary root multiplicities (from the
    automorphic correction by $\phi_{0,1}$) forces all
    higher arities to be nonzero.
\end{itemize}

\noindent
The fibration thus produces a \textbf{shadow depth upgrade}:
\[
  \text{class } G \;\;(\text{fiber}) \quad\longrightarrow\quad
  \text{class } M \;\;(\text{total space}).
\]
The mechanism: the base direction $B$ introduces a new
``spectral parameter'' (the elliptic modular parameter
$\sigma$) which couples the fiber shadow to the Jacobi
variable $z$, producing imaginary roots whose multiplicities
grow without bound.

\subsection{The Borcherds lift as shadow completion}

The Borcherds multiplicative lift
$\phi_{0,1}(\tau, z) \mapsto \Delta_5(\tau, z, \sigma)$
is, in the shadow tower language, the passage from the
finite (arity-2) shadow of the fiber to the infinite shadow
tower of the total space.  Concretely:

\begin{enumerate}[label=(\roman*)]
  \item Start with the K3 elliptic genus $\phi_{0,1}$, which
    encodes the $\CY_2$ root multiplicities $f(n,l)$.

  \item The fiber shadow at arity 2 gives
    $\kappa_{\text{fiber}} = \chi_{\mathrm{top}}(K3) = 24$
    (the topological Euler characteristic of K3; note
    $f(0,0) = 20 = h^{1,1}(K3)$, but $\kappa_{\text{fiber}}$
    receives contributions from all Hodge numbers).

  \item The Borcherds product
    $\Phi_5 = q^a y^b p^c \prod (1 - q^n y^l p^m)^{f(nm,l)}$
    is the bar-complex Euler product of $A_X$: each factor
    encodes one arity of the shadow tower, with the product
    $nm$ implementing the depth--arity transfer.

  \item The global $\kappa_{\text{global}} = 5 = \mathrm{wt}(\Delta_5)$
    is NOT $\kappa_{\text{fiber}}/\text{(something)}$; it is
    an independent invariant of the total space, determined by
    the Weyl vector $\rho_3$ of the $\CY_3$ BKM algebra.
\end{enumerate}

\begin{remark}[Fiber $\kappa_{\mathrm{ch}}$ vs global $\kappa_{\mathrm{ch}}$]
The dramatic change $\kappa_{\mathrm{ch}} = 24 \to \kappa_{\mathrm{ch}} = 5$ is not a
``renormalization'' of $\kappa_{\mathrm{ch}}$ by the base.  Rather, the
fiber $\kappa_{\mathrm{ch}}$ and the global $\kappa_{\mathrm{ch}}$ are invariants of
\emph{different} objects ($A_S$ vs $A_X$), and the map
$A_S \mapsto A_X$ (incorporating the elliptic fibration) is
not $\kappa_{\mathrm{ch}}$-linear.  The global $\kappa_{\mathrm{ch}} = 5$ is determined
by the automorphic weight of the Siegel modular form, which
depends on the full root system of $\mathfrak{g}_{\Delta_5}$,
not on the fiber alone.
\end{remark}

\end{document}
